\documentclass[12pt]{article}
\usepackage[utf8]{inputenc}
\usepackage[T1]{fontenc}
\usepackage{hyperref}
\title{Yanom\'ami Language Resources for K-12 Curriculum Development}
\author{Compiled by Claude Sonnet 5 (claude-sonnet-5) for Andr\'as Kornai}
\date{}
\begin{document}
\maketitle

\section{Introduction}
Yanom\'am (ISO 639-3 wca) is one of the closely related Yanomaman languages spoken in northern Brazil. Most of its speakers are reported to be monolingual. Searches found a few reference works and community-organisation sources, but no primer or school curriculum.

\section{Primary Educational Resources}
No primer or textbook was found. The dictionary below was reported as intended to help Yanomami people learn Portuguese.

\section{Linguistic and Reference Materials}
\begin{itemize}
\item Portuguese-language dictionary hosted by the Instituto Socioambiental (ISA), reported as finalized in 2011 and covering 666 Portuguese verbs: \url{https://acervo.socioambiental.org/sites/default/files/publications/yal00020.pdf}
\item Yanomami--Spanish / Spanish--Yanomami illustrated dictionary by Cristina Martins Fargetti, reviewed in \emph{Mana} (2008). Repository record: \url{https://acervodigital.unesp.br/handle/11449/6410}
\item \emph{L\'inguas Yanomami no Brasil}, from the ISA and Hutukara linguistic diversity project, which identified a sixth language, Y\~anoma: \url{https://www.nationalgeographicbrasil.com/cultura/2023/05/que-lingua-os-yanomamis-falam}
\item ISA article on the Yanomami languages: \url{https://site-antigo.socioambiental.org/en/print/6890}
\item Wikipedia on Yanom\'am: \url{https://en.wikipedia.org/wiki/Yanom%C3%A1m_language}
\item Wikipedia on the Yanomaman family: \url{https://en.wikipedia.org/wiki/Yanomaman_languages}
\end{itemize}

\section{Religious and Community Resources}
\begin{itemize}
\item Hutukara Associa\c{c}\~ao Yanomami, the Yanomami community organisation. A document on Hutukara was seen in the results: \url{https://ea.fflch.usp.br/sites/ea.fflch.usp.br/files/inline-files/Hutukara%20EN.pdf}
\end{itemize}

\section{Digital and Audio Resources}
None found in the search results.

\section{Curriculum Development Context}
The family is large and divided into several languages and at least sixteen dialect variations. A curriculum must therefore be tied to the specific variety taught, and the ISA and Hutukara partnership is the clearest route. The school status of Yanom\'am is unconfirmed here, and no primers were found.
\end{document}